% =====================================================================
%  Chương 5 — Mục 5 của đề: so sánh và đánh giá
%  Mã nguồn: a05/metrics.py · Notebook: 03, 04, 05 (số liệu), 06 (tổng hợp)
% =====================================================================
\chapter{So sánh và đánh giá}
\label{ch:sosanh}

\section{Tiêu chí đánh giá}

Accuracy một mình không đủ để so sánh các mô hình: nó bỏ qua giá phải trả, và gần như vô nghĩa khi dữ liệu mất cân bằng
như Diabetes. Chương này so sánh bốn mô hình theo năm nhóm tiêu chí:

\begin{enumerate}
  \item \textbf{Chất lượng} trên tập test: accuracy và macro-F1 với ảnh; thêm top-5 với CIFAR-100 (100 lớp, nhiều lớp
        gần giống nhau); precision, recall, F1, ROC-AUC và PR-AUC với Diabetes.
  \item \textbf{Chi phí}: số tham số, số phép nhân-cộng (MACs) cho một mẫu, thời gian một epoch, thời gian suy luận
        một mẫu.
  \item \textbf{Hội tụ}: đường loss và accuracy theo epoch, epoch có val loss nhỏ nhất.
  \item \textbf{Lỗi}: ma trận nhầm lẫn của mô hình tốt nhất.
  \item \textbf{Biểu diễn}: bộ lọc lớp đầu và feature map của cùng một ảnh qua bốn mô hình.
\end{enumerate}

Macro-F1 là trung bình F1 của từng lớp, nên một lớp bị bỏ quên sẽ kéo điểm này xuống dù accuracy vẫn cao. Top-$k$
accuracy tính tỉ lệ mẫu có nhãn đúng nằm trong $k$ lớp có điểm cao nhất.

\maNguon{metrics__topk_accuracy}{Top-$k$ accuracy}
\maNguon{metrics__classification_report}{Ba chỉ số chất lượng cho ảnh}

Số MACs chỉ tính các tầng tích chập và nối đầy đủ; BatchNorm, ReLU, pooling và phép cộng đường tắt nhỏ hơn vài bậc nên
được bỏ qua. Thời gian suy luận được đo trên lô 256 ảnh, có đồng bộ GPU trước và sau khi đo, rồi chia cho số ảnh.

\maNguon{metrics__count_macs}{Đếm MACs bằng forward hook}
\maNguon{metrics__time_inference}{Đo thời gian suy luận có đồng bộ GPU}

Vì mỗi mô hình chỉ chạy một hạt giống, chênh lệch nhỏ giữa hai mô hình có thể chỉ là nhiễu của lần chạy đó. Chương này
quy ước: chênh lệch dưới 0,5 điểm phần trăm được ghi là \emph{nằm trong mức dao động của một lần chạy} và không dùng để
kết luận mô hình nào tốt hơn.

\section{CIFAR-10}

\begin{table}[H]
  \centering
  \caption{Kết quả trên tập test CIFAR-10 (\SL{cifar10/n_test} ảnh). In đậm: accuracy cao nhất}
  \label{tab:kqcifar10}
  \small\input{generated/tab_cifar10}
\end{table}

\paragraph{Chất lượng.} Bước lớn nhất là từ M0 sang M1: accuracy tăng \SL{compare/delta/cifar10/basic_vgg} điểm, từ
\SL{cifar10/models/basic/acc}\,\% lên \SL{cifar10/models/vgg/acc}\,\%. Hai bước sau không cho lợi ích nào: thêm đường tắt
đổi \SL{compare/delta/cifar10/vgg_resnet} điểm, trong mức dao động. Thêm SE đổi \SL{compare/delta/cifar10/resnet_seresnet}
điểm, nhỉnh hơn ngưỡng 0,5 một chút và theo chiều xấu đi. Tuy vậy không nên đọc đó là tác hại của SE: trên tập validation,
thứ tự của ba mô hình ở epoch cuối lại ngược với trên tập test, và SE-ResNet đứng đầu (\SL{compare/final/cifar10/vgg/val_acc}, \SL{compare/final/cifar10/resnet/val_acc} và
\SL{compare/final/cifar10/seresnet/val_acc}\,\% cho M1, M2, M3). Macro-F1 gần như trùng accuracy vì mười lớp cân bằng.

\paragraph{Vì sao đường tắt không giúp gì ở đây.} ResNet giải quyết bài toán \emph{tối ưu mạng rất sâu}: khi chồng hàng
chục tầng, gradient khó đi về các tầng đầu. M1 chỉ có 12 tầng tích chập, lại có BatchNorm, nên bài toán đó chưa xuất hiện.
M1 đã đạt accuracy \SL{compare/final/cifar10/vgg/train_acc}\,\% trên tập train ở epoch cuối, tức là tối ưu không còn là
chỗ nghẽn. Đường tắt vì vậy không có gì để sửa. Kết quả này không mâu thuẫn với bài báo ResNet; nó cho thấy lợi ích của
một cơ chế phụ thuộc vào việc mô hình có đang gặp đúng vấn đề mà cơ chế đó giải quyết hay không.

\paragraph{Hội tụ.} Hình~\ref{fig:hoccifar10} cho thấy M0 hội tụ chậm và dừng ở accuracy train
\SL{compare/final/cifar10/basic/train_acc}\,\%, trong khi ba mô hình còn lại gần như thuộc lòng tập train. Khoảng cách
train--val của M0 là \SL{compare/final/cifar10/basic/gap} điểm, của M1 là \SL{compare/final/cifar10/vgg/gap} điểm. M0 không
thiếu dữ liệu mà thiếu sức chứa: nó chưa khớp hết tập train. Cả bốn mô hình đều có epoch tốt nhất ở cuối lịch huấn luyện,
nhờ pha giảm tốc độ học của OneCycle.

\begin{figure}[H]
  \centering
  \hinh{f5_duong_hoc_cifar10}
  \caption{Đường học trên CIFAR-10}
  \label{fig:hoccifar10}
\end{figure}

\section{CIFAR-100}

\begin{table}[H]
  \centering
  \caption{Kết quả trên tập test CIFAR-100 (\SL{cifar100/n_test} ảnh, 100 lớp)}
  \label{tab:kqcifar100}
  \small\input{generated/tab_cifar100}
\end{table}

\paragraph{Chất lượng.} Cùng bốn kiến trúc và cùng cấu hình, chỉ tăng số lớp lên mười lần, accuracy của mọi mô hình giảm
mạnh so với CIFAR-10 (Bảng~\ref{tab:kqcifar10} và~\ref{tab:kqcifar100}). Bước M0$\to$M1 lớn hơn hẳn: \SL{compare/delta/cifar100/basic_vgg} điểm, từ
\SL{cifar100/models/basic/acc}\,\% lên \SL{cifar100/models/vgg/acc}\,\%. Khác với CIFAR-10, đường tắt lần này có tác dụng
đo được: M2 hơn M1 \SL{compare/delta/cifar100/vgg_resnet} điểm, vượt ngưỡng 0,5 điểm, và top-5 cũng tăng theo (từ
\SL{cifar100/models/vgg/top5} lên \SL{cifar100/models/resnet/top5}\,\%). SE vẫn chỉ thêm
\SL{compare/delta/cifar100/resnet_seresnet} điểm, trong mức dao động.

Một cách đọc hợp lý là: bài toán 100 lớp khó hơn nên M1 không còn gần như thuộc lòng tập train như ở CIFAR-10
(\SL{compare/final/cifar100/vgg/train_acc}\,\% so với \SL{compare/final/cifar10/vgg/train_acc}\,\%), và đường tắt giúp
tối ưu tốt hơn một chút. Tuy vậy, với một lần chạy duy nhất, chênh lệch một điểm chỉ đủ để nói ``có dấu hiệu'', chưa đủ để
kết luận chắc chắn.

\paragraph{Top-5.} Top-5 của mô hình tốt nhất là \SL{cifar100/models/seresnet/top5}\,\%, cao hơn nhiều so với top-1.
Phần lớn lỗi top-1 là nhầm sang một lớp gần giống, và lớp đúng vẫn nằm trong vài lựa chọn đầu. Mục~\ref{sec:loi} sẽ cho thấy
các lớp gần giống này thường thuộc cùng một siêu lớp.

\paragraph{Hội tụ.} Khoảng cách train--val trên CIFAR-100 rất lớn: \SL{compare/final/cifar100/vgg/gap} điểm với M1 và
\SL{compare/final/cifar100/basic/gap} điểm với M0 (Hình~\ref{fig:hoccifar100}). Mỗi lớp chỉ có \SL{data/cifar100/per_class_max} ảnh gốc, nên các
mô hình có BN nhớ được tập train nhưng khái quát kém hơn nhiều so với CIFAR-10. Đây là quá khớp thật, khác với trường hợp
thiếu sức chứa của M0.

\begin{figure}[H]
  \centering
  \hinh{f5_duong_hoc_cifar100}
  \caption{Đường học trên CIFAR-100}
  \label{fig:hoccifar100}
\end{figure}

\section{Diabetes}

\begin{table}[H]
  \centering
  \caption{Kết quả trên tập test Diabetes (\SL{diabetes/n_test} bệnh nhân). In đậm: giá trị cao nhất mỗi cột}
  \label{tab:kqdiabetes}
  \scriptsize\input{generated/tab_diabetes}
\end{table}

\paragraph{Mốc so sánh.} Đoán mọi bệnh nhân đều khoẻ đạt accuracy \SL{diabetes/majority_acc}\,\% nhưng recall bằng 0. Accuracy của bốn
mô hình từ \SL{diabetes/models/basic/acc} đến \SL{diabetes/models/seresnet/acc}\,\%, chỉ hơn mốc này vài điểm. Vì vậy phần so sánh dưới đây dựa vào recall, F1 và
hai diện tích dưới đường cong.

\paragraph{Chất lượng.} Bốn mô hình gần như không phân biệt được. ROC-AUC chỉ dao động từ \SL{diabetes/models/vgg/roc_auc}
đến \SL{diabetes/models/basic/roc_auc}, và PR-AUC từ \SL{diabetes/models/vgg/pr_auc} đến \SL{diabetes/models/seresnet/pr_auc}.
Hai chỉ số này đo khả năng \emph{xếp hạng} bệnh nhân theo nguy cơ và không phụ thuộc ngưỡng, nên có thể nói bốn mô hình xếp
hạng tốt như nhau. F1 của SE-ResNet cao hơn (\SL{diabetes/models/seresnet/f1}\,\% so với \SL{diabetes/models/basic/f1}\,\% của M0),
nhưng Bảng~\ref{tab:kqdiabetes} cho thấy phần chênh này đến từ \emph{điểm cắt}: precision của nó cao nhất
(\SL{diabetes/models/seresnet/precision}\,\%) còn recall thấp nhất (\SL{diabetes/models/seresnet/recall}\,\%). Nó không phải
một mô hình tốt hơn, chỉ là một điểm khác trên cùng một đường ROC (Hình~\ref{fig:rocpr}).

\begin{figure}[H]
  \centering
  \hinh{f5_roc_pr}
  \caption{Đường ROC và Precision--Recall của bốn mô hình trên tập test Diabetes. Đường chấm trên hình PR là tỉ lệ dương, mốc của một bộ phân loại ngẫu nhiên}
  \label{fig:rocpr}
\end{figure}

\paragraph{Vì sao kiến trúc không giúp gì ở đây.} Cả ba cơ chế đều được thiết kế cho dữ liệu có cấu trúc không gian: bộ lọc
$3\times3$ khai thác quan hệ giữa các điểm ảnh kề nhau, và một đặc trưng dịch đi vài ô vẫn là cùng một đặc trưng. Một hàng
dữ liệu Diabetes chỉ có \SL{data/diabetes/n_features_encoded} cột, và hai cột kề nhau (ví dụ \texttt{bmi} và
\texttt{HbA1c\_level}) không có quan hệ lân cận nào. Chương~\ref{ch:dulieu} cũng cho thấy phần lớn ca dương tính đã được
hai ngưỡng xét nghiệm tách ra. Phần khó còn lại nằm ở những bệnh nhân có chỉ số chồng lấn với người khoẻ, và thêm tầng hay
thêm cơ chế không tạo ra thông tin mới để phân biệt họ.

\paragraph{Hội tụ.} Các mô hình có BN đạt val loss nhỏ nhất rất sớm (epoch \SL{diabetes/models/vgg/best_epoch} với M1,
\SL{diabetes/models/resnet/best_epoch} với M2) rồi bị dừng sớm; M0 cần đến epoch \SL{diabetes/models/basic/best_epoch}
(Hình~\ref{fig:hocdiabetes}). Khoảng cách train--val ở epoch cuối nhỏ (\SL{compare/final/diabetes/basic/gap} điểm với M0): với \SL{data/diabetes/n_train} dòng
train và nhiều nhất \SL{diabetes/models/seresnet/params} tham số, không mô hình nào quá khớp.

\begin{figure}[H]
  \centering
  \hinh{f5_duong_hoc_diabetes}
  \caption{Đường học trên Diabetes; accuracy tính tại ngưỡng 0,5 của mô hình có \texttt{pos\_weight}}
  \label{fig:hocdiabetes}
\end{figure}

\paragraph{Ma trận nhầm lẫn.} Tại ngưỡng chọn trên validation, SE-ResNet bỏ sót \SL{diabetes/models/seresnet/confusion/1/0}
ca dương tính và báo nhầm \SL{diabetes/models/seresnet/confusion/0/1} người khoẻ; M0 bỏ sót
\SL{diabetes/models/basic/confusion/1/0} ca và báo nhầm \SL{diabetes/models/basic/confusion/0/1} người. Nếu mô hình được dùng
để sàng lọc, bỏ sót một ca bệnh đắt hơn nhiều so với một lần xét nghiệm thừa, và khi đó nên hạ ngưỡng để tăng recall thay vì
chọn ngưỡng theo F1.

\section{Chất lượng và chi phí}
\label{sec:chiphi}

Hình~\ref{fig:cot} đặt chất lượng cạnh thời gian huấn luyện. Bước M0$\to$M1 tăng số tham số từ
\SL{cifar10/models/basic/paramsM} lên \SL{cifar10/models/vgg/paramsM} triệu và số MACs mỗi ảnh từ
\SL{cifar10/models/basic/macsM} lên \SL{cifar10/models/vgg/macsM} triệu. Phần lớn MACs tăng thêm đến từ việc các tầng
đầu của M1 chạy ở độ phân giải $32\times32$ với 64 kênh. Hai bước sau gần như miễn phí về tham số: đường tắt thêm
\SL{compare/delta/cifar10/vgg_resnet_params} tham số (các conv $1\times1$ chiếu) và SE thêm
\SL{compare/delta/cifar10/resnet_seresnet_params}. Tuy vậy thời gian mỗi epoch vẫn tăng, từ \SL{cifar10/models/vgg/epoch_s}
lên \SL{cifar10/models/resnet/epoch_s} giây, vì phép cộng đường tắt và SE là các thao tác đọc ghi bộ nhớ, không được tính
vào MACs. Trên Diabetes, M0 huấn luyện một epoch trong \SL{diabetes/models/basic/epoch_s} giây, còn SE-ResNet cần
\SL{diabetes/models/seresnet/epoch_s} giây, cho cùng một chất lượng.

\begin{figure}[H]
  \centering
  \hinh{f5_cot}
  \caption{Chất lượng (hàng trên) và thời gian mỗi epoch (hàng dưới) của M0--M3 trên ba tập}
  \label{fig:cot}
\end{figure}

Hình~\ref{fig:scatter} vẽ accuracy theo chi phí. Ba mô hình M1--M3 nằm chồng lên nhau ở cùng mức chi phí. Trên cả hai tập ảnh,
đường nối M0 với cụm này dốc lên rõ rệt: số tham số tăng từ \SL{cifar10/models/basic/paramsM} lên \SL{cifar10/models/vgg/paramsM} triệu đổi được \SL{compare/delta/cifar10/basic_vgg}
điểm trên CIFAR-10 và \SL{compare/delta/cifar100/basic_vgg} điểm trên CIFAR-100.

\begin{figure}[H]
  \centering
  \hinh{f5_scatter}
  \caption{Accuracy trên tập test theo số tham số (trái) và theo MACs mỗi ảnh, trục log (phải)}
  \label{fig:scatter}
\end{figure}

\section{Phân tích lỗi}
\label{sec:loi}

Hình~\ref{fig:confc10} là ma trận nhầm lẫn của \SL{compare/confusion/cifar10/model}, mô hình tốt nhất trên CIFAR-10. Lớp
khó nhất là \SL{compare/confusion/cifar10/worst} (recall \SL{compare/confusion/cifar10/diag_min}\,\%), lớp dễ nhất là
\SL{compare/confusion/cifar10/best} (\SL{compare/confusion/cifar10/diag_max}\,\%). Cặp nhầm nhiều nhất là
\SL{compare/confusion/cifar10/top_pair/0} bị đoán thành \SL{compare/confusion/cifar10/top_pair/1}
(\SL{compare/confusion/cifar10/top_pair_pct}\,\% số ảnh \SL{compare/confusion/cifar10/top_pair/0}). Hai lớp này đều là thú bốn chân
có lông, và ở độ phân giải $32\times32$ chúng khác nhau chủ yếu ở hình dạng tai và mõm. Lỗi của mô hình đi theo sự giống
nhau về hình ảnh chứ không rải ngẫu nhiên.

\begin{figure}[H]
  \centering
  \hinh{f5_conf_cifar10}
  \caption{Ma trận nhầm lẫn CIFAR-10 của mô hình tốt nhất, \% theo hàng}
  \label{fig:confc10}
\end{figure}

Với CIFAR-100, ma trận $100\times100$ quá dày để đọc, nên lớp thật và lớp dự đoán được gộp về 20 siêu lớp
(Hình~\ref{fig:confc100}). Mô hình tốt nhất là \SL{compare/confusion/cifar100/model}. Accuracy ở mức siêu lớp là
\SL{compare/confusion/cifar100/coarse_acc}\,\%, cao hơn hẳn top-1 ở mức lớp mịn (\SL{cifar100/models/seresnet/acc}\,\%). Như vậy
rất nhiều lỗi là nhầm giữa hai lớp mịn \emph{trong cùng một siêu lớp}: dự đoán sai lớp nhưng vẫn đúng nhóm. Cặp siêu lớp
nhầm nhiều nhất là \SL{compare/confusion/cifar100/top_pair/0} sang \SL{compare/confusion/cifar100/top_pair/1}
(\SL{compare/confusion/cifar100/top_pair_pct}\,\%). Siêu lớp khó nhất là \SL{compare/confusion/cifar100/worst}, dễ nhất là
\SL{compare/confusion/cifar100/best}.

\begin{figure}[H]
  \centering
  \hinh{f5_conf_cifar100}
  \caption{Ma trận nhầm lẫn CIFAR-100 gộp theo 20 siêu lớp, \% theo hàng}
  \label{fig:confc100}
\end{figure}

\section{Biểu diễn: bộ lọc lớp đầu và feature map}

Hình~\ref{fig:bieudien} đặt cạnh nhau 16 bộ lọc $3\times3\times3$ đầu tiên của tầng tích chập đầu và sáu feature map mà tầng đó
tạo ra từ cùng một ảnh test (lớp \SL{compare/probe_class}, ảnh ở Hình~\ref{fig:fmapm0}), cho bốn mô hình CIFAR-10. Mỗi bộ lọc
được kéo giãn min--max về $[0, 255]$ để hiển thị như một ảnh màu nhỏ.

\begin{figure}[H]
  \centering
  \hinh{f5_bieu_dien}
  \caption{Bộ lọc lớp đầu và feature map của cùng một ảnh test qua bốn mô hình CIFAR-10}
  \label{fig:bieudien}
\end{figure}

Ở tầng đầu, bốn mô hình học những thứ giống nhau: các bộ lọc là những mảng màu và những cặp sáng--tối theo một hướng, và
feature map làm nổi đường viền con ngựa theo các hướng khác nhau. Một vài kênh của ResNet và SE-ResNet phản ứng với cả vùng
thân ngựa thay vì chỉ với viền. Tầng đầu của bốn mô hình giống nhau như vậy cũng khớp với việc M1, M2, M3 cho accuracy gần như bằng
nhau: khác biệt giữa chúng, nếu có, nằm ở các tầng sâu hơn.

Có một khác biệt đo được: độ lệch chuẩn của trọng số bộ lọc lớp đầu ở M0 là \SL{compare/filters/basic/weight_std}, còn ở
M1, M2, M3 chỉ từ \SL{compare/filters/seresnet/weight_std} đến \SL{compare/filters/vgg/weight_std}. Điều này có lý do: ở M1--M3,
ngay sau mỗi tầng tích chập là BatchNorm, và BatchNorm chuẩn hoá lại đầu ra, nên nhân mọi trọng số của bộ lọc với một hằng số
không làm đổi kết quả. Độ lớn của trọng số vì thế không còn ý nghĩa, và weight decay kéo nó nhỏ dần. M0 không có BatchNorm nên
độ lớn trọng số quyết định trực tiếp độ lớn của feature map.

\section{Tổng hợp: mỗi cơ chế mang lại bao nhiêu}

Bảng~\ref{tab:delta} gom chênh lệch giữa hai mô hình liền kề trên cả ba tập, và Bảng~\ref{tab:tonghop} là toàn bộ 12 lượt
chạy (3 tập $\times$ 4 mô hình).

\begin{table}[H]
  \centering
  \caption{Chênh lệch chất lượng giữa hai mô hình liền kề}
  \label{tab:delta}
  \small\input{generated/tab_delta}
\end{table}

\begin{table}[H]
  \centering
  \caption{Tổng hợp 12 lượt chạy trên tập test}
  \label{tab:tonghop}
  \small\input{generated/tab_summary}
\end{table}

Ba điều rút ra từ hai bảng trên:

\begin{enumerate}
  \item \textbf{Bước M0$\to$M1 quyết định gần như toàn bộ chênh lệch trên ảnh}: \SL{compare/delta/cifar10/basic_vgg} điểm
        trên CIFAR-10 và \SL{compare/delta/cifar100/basic_vgg} điểm trên CIFAR-100. Bước này gộp ba thay đổi (sâu hơn,
        BatchNorm, GAP), và thí nghiệm không tách được phần đóng góp của từng thay đổi. Bằng chứng liên quan duy nhất
        là M0, mô hình không có BatchNorm, đã phân kỳ ở tốc độ học đỉnh khi chưa cắt gradient (Hình~\ref{fig:clip}).
  \item \textbf{Đường tắt chỉ có ích khi bài toán đủ khó}: không đổi gì trên CIFAR-10 (\SL{compare/delta/cifar10/vgg_resnet}
        điểm), tăng \SL{compare/delta/cifar100/vgg_resnet} điểm trên CIFAR-100. Với mạng 12 tầng, lợi ích này nhỏ; nó được
        thiết kế cho mạng sâu hàng chục đến hàng trăm tầng.
  \item \textbf{SE không cho lợi ích nhất quán}: \SL{compare/delta/cifar10/resnet_seresnet} điểm trên CIFAR-10 (validation
        lại xếp nó đứng đầu), \SL{compare/delta/cifar100/resnet_seresnet} điểm trên CIFAR-100. Trên Diabetes, F1 của nó cao hơn
        \SL{compare/delta/diabetes/resnet_seresnet} điểm nhưng ROC-AUC không đổi, tức là khác biệt nằm ở ngưỡng chứ không ở
        mô hình.
\end{enumerate}

Trên dữ liệu bảng, lựa chọn hợp lý là mô hình \emph{rẻ nhất}: M0 cho cùng khả năng xếp hạng với \SL{diabetes/models/basic/params}
tham số, so với \SL{diabetes/models/seresnet/params} của M3. Trên ảnh, lựa chọn hợp lý là M1 hoặc M2: M2 tốt hơn một chút khi bài toán khó hơn, với chi phí tham số gần như bằng M1.
